
\subsection{Relation to current cross-subject EEG validation research}
This study complements, rather than replaces, the development of
spectral filters, classifiers and adaptation methods. Recent
cross-subject EEG work has demonstrated that window- or trial-level
performance estimates can differ considerably from genuine
unseen-participant performance \cite{gil2026subjectvalidation}. A
separate 45-person mental-workload study has shown that feature
representation, model class and the allowance for subject-specific
calibration can also strongly affect generalisation
\cite{abinaya2026wavelet}. Both results caution against evaluating
one model under a more favourable target-information budget than
another. They do not directly answer the within-pipeline question
of whether selecting \emph{the same} scalar regulariser by
covariance predictive likelihood or by source-person
discrimination changes which EEG representation appears better.

We deliberately keep the comparison narrower than studies introducing
new classifiers. The use of target-person unlabeled moments,
target-conditional geometry, entropy minimisation or a few labeled
calibration trials may benefit model adaptation, but changes the
information contract. A Z0 result cannot be directly ranked against
one of these adapted classifiers as if the procedures used the same
information. In this article, Ledoit--Wolf, OAS and source
covariance-NLL are estimator/objective controls; Riemannian MDM,
tangent logistic regression and a separate CSP6 classifier
are architectural reference points. Comparing those latter
representations does not prove a new state-of-the-art BCI decoder.

The essential methodological contrast is between \emph{estimation
quality of a covariance operator} and \emph{out-of-person
discrimination of a derived decision direction}. The mathematical
non-equivalence is familiar from regularised discriminant analysis
\cite{friedman1989}; our empirical contribution is its prospective,
participant-level consequence, including a spectral representation
hypothesis that the new results contradict. This framing prevents
a conventional parameter-selection strategy from being presented
as a newly invented learning algorithm.
